% GENERATED FILE. Edit canonical lessons, then run npm run generate:books.
% canonical-source-hash: fnv1a64:2358812752f78736
% canonical-lessons: MW-C183-karchhi, MW-C183-tokri, MW-C183-angochho, MW-C183-kajal, MW-C183-besan

\chapter{Small Things Around the House}
\label{ch:mw-small-things-around-the-house}
\hlchaptermodality{183}

\begin{chapteropening}
\textbf{By the end of this chapter:} \emph{I can name small things I use around the house.}

Complete the last lesson of chapter 183: I can name small things I use around the house.
\end{chapteropening}

\section[\texorpdfstring{\mw{कड़छी} (kaṛchhī)}{कड़छी (kaṛchhī)}]{\mw{कड़छी} (kaṛchhī) — a ladle}

\label{lesson:MW-C183-karchhi}



\begin{quote}
Before the new one: say the Marwadi for to temper, then the Marwadi for to serve food.
\end{quote}

\subsection*{You'll want to know: \mw{कड़छी}}
\textbf{\mw{कड़छी}} — \emph{kaṛchhī} — \textquotedblleft{}a ladle\textquotedblright{}.

\begin{quote}
\mw{कड़छी} \mw{कड़ाही} \mw{मांय} \mw{है।}
\end{quote}

\begin{quote}
\emph{kaṛchhī kaṛāhī mā̃y hai.}
\end{quote}

\begin{quote}
The ladle is in the wok.
\end{quote}

A feminine noun.

\begin{cousinweb}
\textbf{\mw{कड़छी}} (\emph{kaṛchhī}) comes from a reconstructed \emph{*kaṭacchu}, \textquotedblleft{}a ladle\textquotedblright{}; CDIAL lists Hindi and Gujarati \emph{kaṛchī}, \textquotedblleft{}an iron ladle\textquotedblright{} (CDIAL 2633). CDIAL gives no Marwari form in this entry.
\end{cousinweb}

\begin{grammarlens}[title={What you've built}]
One more word for this chapter.
\end{grammarlens}

\subsection*{Your turn}
\begin{itemize}
\raggedright
  \item \emph{Say it:} \emph{kaṛchhī}
  \item \emph{Say it:} \emph{kaṛchhī}, once more
  \item \emph{Say it:} \emph{kaṛchhī kaṛāhī mā̃y hai.} — The ladle is in the wok.
  \item \emph{Recall:} say the Marwadi for to temper, then the Marwadi for to serve food, then say \emph{kaṛchhī} again
\end{itemize}

\begin{tcolorbox}[breakable,colback=teal!4,colframe=teal!35!black,title={Before you move on}]
What does \mw{कड़छी} mean? (\textquotedblleft{}A ladle\textquotedblright{}.) Say it once more.
\end{tcolorbox}
\section[\texorpdfstring{\mw{टोकरी} (ṭokrī)}{टोकरी (ṭokrī)}]{\mw{टोकरी} (ṭokrī) — a basket}

\label{lesson:MW-C183-tokri}



\begin{quote}
Before the new one: say the Marwadi for to serve food, then the Marwadi for a ladle.
\end{quote}

\subsection*{You'll want to know: \mw{टोकरी}}
\textbf{\mw{टोकरी}} — \emph{ṭokrī} — \textquotedblleft{}a basket\textquotedblright{}.

\begin{quote}
\mw{टोकरी} \mw{मांय} \mw{सबजी} \mw{है।}
\end{quote}

\begin{quote}
\emph{ṭokrī mā̃y sabjī hai.}
\end{quote}

\begin{quote}
There are vegetables in the basket.
\end{quote}

A feminine noun.

\begin{cousinweb}
\textbf{\mw{टोकरी}} (\emph{ṭokrī}) comes from a reconstructed \emph{*ṭōkka}, \textquotedblleft{}a basket\textquotedblright{}; CDIAL lists Hindi \emph{ṭokrī} (CDIAL 5477). CDIAL gives no Marwari form in this entry.
\end{cousinweb}

\begin{grammarlens}[title={What you've built}]
One more word for this chapter.
\end{grammarlens}

\subsection*{Your turn}
\begin{itemize}
\raggedright
  \item \emph{Say it:} \emph{ṭokrī}
  \item \emph{Say it:} \emph{ṭokrī}, once more
  \item \emph{Say it:} \emph{ṭokrī mā̃y sabjī hai.} — There are vegetables in the basket.
  \item \emph{Recall:} say the Marwadi for to serve food, then the Marwadi for a ladle, then say \emph{ṭokrī} again
\end{itemize}

\begin{tcolorbox}[breakable,colback=teal!4,colframe=teal!35!black,title={Before you move on}]
What does \mw{टोकरी} mean? (\textquotedblleft{}A basket\textquotedblright{}.) Say it once more.
\end{tcolorbox}
\section[\texorpdfstring{\mw{अंगोछो} (aṅgochho)}{अंगोछो (aṅgochho)}]{\mw{अंगोछो} (aṅgochho) — a shoulder cloth}

\label{lesson:MW-C183-angochho}



\begin{quote}
Before the new one: say the Marwadi for a ladle, then the Marwadi for a basket.
\end{quote}

\subsection*{You'll want to know: \mw{अंगोछो}}
\textbf{\mw{अंगोछो}} — \emph{aṅgochho} — \textquotedblleft{}a shoulder cloth\textquotedblright{}.

\begin{quote}
\mw{म्हारो} \mw{अंगोछो} \mw{गीलो} \mw{है।}
\end{quote}

\begin{quote}
\emph{mhāro aṅgochho gīlo hai.}
\end{quote}

\begin{quote}
My shoulder cloth is wet.
\end{quote}

A masculine noun.

\begin{cousinweb}
\textbf{\mw{अंगोछो}} (\emph{aṅgochho}) is inherited from Sanskrit \emph{aṅgoñcha}, \textquotedblleft{}a towel\textquotedblright{}; CDIAL lists Hindi \emph{ãgochā} and Sindhi \emph{aṅgocho} (CDIAL 139). CDIAL gives no Marwari form in this entry.
\end{cousinweb}

\begin{grammarlens}[title={What you've built}]
One more word for this chapter.
\end{grammarlens}

\subsection*{Your turn}
\begin{itemize}
\raggedright
  \item \emph{Say it:} \emph{aṅgochho}
  \item \emph{Say it:} \emph{aṅgochho}, once more
  \item \emph{Say it:} \emph{mhāro aṅgochho gīlo hai.} — My shoulder cloth is wet.
  \item \emph{Recall:} say the Marwadi for a ladle, then the Marwadi for a basket, then say \emph{aṅgochho} again
\end{itemize}

\begin{tcolorbox}[breakable,colback=teal!4,colframe=teal!35!black,title={Before you move on}]
What does \mw{अंगोछो} mean? (\textquotedblleft{}A shoulder cloth\textquotedblright{}.) Say it once more.
\end{tcolorbox}
\section[\texorpdfstring{\mw{काजळ} (kājaḷ)}{काजळ (kājaḷ)}]{\mw{काजळ} (kājaḷ) — kohl}

\label{lesson:MW-C183-kajal}



\begin{quote}
Before the new one: say the Marwadi for a basket, then the Marwadi for a shoulder cloth.
\end{quote}

\subsection*{You'll want to know: \mw{काजळ}}
\textbf{\mw{काजळ}} — \emph{kājaḷ} — \textquotedblleft{}kohl\textquotedblright{}.

\begin{quote}
\mw{छोरी} \mw{आंख} \mw{मांय} \mw{काजळ} \mw{लगावै} \mw{है।}
\end{quote}

\begin{quote}
\emph{chhorī ā̃kh mā̃y kājaḷ lagāvai hai.}
\end{quote}

\begin{quote}
The girl puts kohl in her eyes.
\end{quote}

A masculine noun.

\begin{cousinweb}
\textbf{\mw{काजळ}} (\emph{kājaḷ}) is inherited from Sanskrit \emph{kajjala}, \textquotedblleft{}lamp-black\textquotedblright{}, a word CDIAL thinks may come from Dravidian or Munda; it lists Gujarati \emph{kājaḷ} and Hindi \emph{kājal} (CDIAL 2622). CDIAL gives no Marwari form in this entry.
\end{cousinweb}

\begin{grammarlens}[title={What you've built}]
One more word for this chapter.
\end{grammarlens}

\subsection*{Your turn}
\begin{itemize}
\raggedright
  \item \emph{Say it:} \emph{kājaḷ}
  \item \emph{Say it:} \emph{kājaḷ}, once more
  \item \emph{Say it:} \emph{chhorī ā̃kh mā̃y kājaḷ lagāvai hai.} — The girl puts kohl in her eyes.
  \item \emph{Recall:} say the Marwadi for a basket, then the Marwadi for a shoulder cloth, then say \emph{kājaḷ} again
\end{itemize}

\begin{tcolorbox}[breakable,colback=teal!4,colframe=teal!35!black,title={Before you move on}]
What does \mw{काजळ} mean? (\textquotedblleft{}Kohl\textquotedblright{}.) Say it once more.
\end{tcolorbox}
\section[\texorpdfstring{\mw{बेसण} (besaṇ)}{बेसण (besaṇ)}]{\mw{बेसण} (besaṇ) — gram flour}

\label{lesson:MW-C183-besan}



\begin{quote}
Before the new one: say the Marwadi for a shoulder cloth, then the Marwadi for kohl.
\end{quote}

\subsection*{You'll want to know: \mw{बेसण}}
\textbf{\mw{बेसण}} — \emph{besaṇ} — \textquotedblleft{}gram flour\textquotedblright{}.

\begin{quote}
\mw{बेसण} \mw{पीळो} \mw{है।}
\end{quote}

\begin{quote}
\emph{besaṇ pīḷo hai.}
\end{quote}

\begin{quote}
Gram flour is yellow.
\end{quote}

A masculine noun.

\begin{cousinweb}
\textbf{\mw{बेसण}} (\emph{besaṇ}) is inherited from Sanskrit \emph{vesana}, a kind of flour; CDIAL lists Gujarati \emph{vesaṇ} and Hindi \emph{besan}, \textquotedblleft{}gram flour\textquotedblright{} (CDIAL 12133). CDIAL gives no Marwari form in this entry.
\end{cousinweb}

\begin{grammarlens}[title={What you've built}]
One more word for this chapter.
\end{grammarlens}

\subsection*{Your turn}
\begin{itemize}
\raggedright
  \item \emph{Say it:} \emph{besaṇ}
  \item \emph{Say it:} \emph{besaṇ}, once more
  \item \emph{Say it:} \emph{besaṇ pīḷo hai.} — Gram flour is yellow.
  \item \emph{Recall:} say the Marwadi for a shoulder cloth, then the Marwadi for kohl, then say \emph{besaṇ} again
\end{itemize}

\begin{tcolorbox}[breakable,colback=teal!4,colframe=teal!35!black,title={Before you move on}]
What does \mw{बेसण} mean? (\textquotedblleft{}Gram flour\textquotedblright{}.) Say it once more.
\end{tcolorbox}
